\documentclass[11pt,a4paper]{article}
\usepackage{shared/luxcover}
\usepackage[utf8]{inputenc}
\usepackage[T1]{fontenc}
\usepackage{amsmath,amssymb,amsfonts,amsthm}
\usepackage{graphicx}
\usepackage{hyperref}
\usepackage{algorithm}
\usepackage{algpseudocode}
\usepackage{booktabs}
\usepackage{listings}
\input{shared/lstlang}
\usepackage{xcolor}
\usepackage{enumitem}
\usepackage{float}

\newtheorem{theorem}{Theorem}
\newtheorem{lemma}[theorem]{Lemma}
\newtheorem{definition}{Definition}

\lstset{
    basicstyle=\small\ttfamily,
    keywordstyle=\color{blue}\bfseries,
    commentstyle=\color{gray},
    stringstyle=\color{red},
    breaklines=true,
    frame=single,
    numbers=left,
    numberstyle=\tiny\color{gray}
}

\hypersetup{
    colorlinks=true,
    linkcolor=blue,
    citecolor=blue,
    urlcolor=blue
}

\title{AES-256-GCM Precompile:\\Authenticated Encryption on the EVM}
\author{
Zach Kelling\thanks{Corresponding author: zach@lux.network}\\
\textit{Lux Industries, Inc.}\\
\texttt{research@lux.network}
}
\date{April 2026}

\begin{document}

\luxcoverpage

\begin{abstract}
We specify the AES-256-GCM precompile deployed on the Lux EVM, providing
authenticated encryption and decryption with associated data (AEAD) as specified
in NIST SP~800-38D. The precompile accepts a 256-bit key, a 96-bit nonce,
plaintext (or ciphertext with tag), and optional associated data, returning the
authenticated ciphertext (or verified plaintext). On-chain symmetric encryption
enables encrypted state variables readable only by designated parties, sealed-bid
auctions, FHE bootstrap key distribution, and confidential governance votes.
Gas cost scales linearly: 3\,000 base plus 6 gas per byte of plaintext/ciphertext,
reflecting the hardware-accelerated AES-NI throughput of modern processors.
\end{abstract}

\section{Introduction}

AES-256-GCM~\cite{sp800-38d} is the industry-standard authenticated encryption
algorithm, providing both confidentiality and integrity in a single operation.
It is mandatory in TLS~1.3, used in IPsec, and supported by all modern CPUs
via the AES-NI instruction set.

Providing AES-256-GCM as an EVM precompile is motivated by:
\begin{enumerate}[label=(\roman*)]
  \item \textbf{Encrypted state}: Contracts can store ciphertext that only
    authorized parties (holding the decryption key) can read.
  \item \textbf{Sealed-bid auctions}: Bids are AES-encrypted on-chain; the
    auctioneer reveals the key at auction close.
  \item \textbf{FHE key material}: Bootstrap keys for FHE schemes are large
    ($\sim$50MB); AES encryption enables efficient transfer and storage.
  \item \textbf{Confidential parameters}: Oracle data can be encrypted for
    specific consumers.
\end{enumerate}

Implementing AES-256-GCM in Solidity requires $\sim$10 rounds of SubBytes,
ShiftRows, MixColumns, and AddRoundKey per 16-byte block, costing $\sim$200\,000
gas per kilobyte. The precompile reduces this to $\sim$6\,000 gas per kilobyte.

\section{Specification}

\begin{definition}[AES-256-GCM]
Parameters:
\begin{itemize}
  \item Key $K$: 32 bytes (256 bits)
  \item Nonce $N$: 12 bytes (96 bits)
  \item Plaintext $P$: 0 to 65\,536 bytes
  \item Associated data $A$: 0 to 1\,024 bytes
  \item Tag $T$: 16 bytes (128-bit authentication tag)
\end{itemize}
\end{definition}

\subsection{Address and Operations}

Address: \texttt{0x0200\ldots25}.

\begin{table}[H]
\centering
\begin{tabular}{llr}
\toprule
\textbf{Selector} & \textbf{Operation} & \textbf{Gas} \\
\midrule
\texttt{0x01} & Encrypt & $3\,000 + 6 \cdot |P|$ \\
\texttt{0x02} & Decrypt & $3\,000 + 6 \cdot |C|$ \\
\bottomrule
\end{tabular}
\end{table}

\subsection{Encrypt}

\textbf{Input}: Selector (1 byte), key $K$ (32 bytes), nonce $N$ (12 bytes),
AAD length (2 bytes, big-endian), AAD ($|A|$ bytes), plaintext ($|P|$ bytes).

\textbf{Output}: Ciphertext ($|P|$ bytes) $\|$ tag (16 bytes).

\subsection{Decrypt}

\textbf{Input}: Selector (1 byte), key $K$ (32 bytes), nonce $N$ (12 bytes),
AAD length (2 bytes), AAD ($|A|$ bytes), ciphertext ($|C|$ bytes) $\|$ tag (16 bytes).

\textbf{Output}: Plaintext ($|C| - 16$ bytes) if tag verifies, error otherwise.

\textbf{Authentication failure}: If the tag does not verify, the precompile returns
an error and consumes all supplied gas. It does \emph{not} return the unauthenticated
plaintext.

\section{Security Analysis}

\subsection{Nonce Reuse}

GCM is catastrophically vulnerable to nonce reuse: reusing $(K, N)$ with different
plaintexts reveals the XOR of plaintexts and enables tag forgery via polynomial
GCD. The precompile does not enforce nonce uniqueness (it cannot---nonce management
is a caller responsibility). However, the gas cost model does not incentivize
nonce reuse since each call costs the same regardless.

\begin{theorem}[GCM security bound]
AES-256-GCM provides IND-CPA and INT-CTXT security with advantage bounded by:
\[
  \text{Adv} \leq \frac{q^2}{2^{128}} + \frac{q \cdot \sigma}{2^{128}} + \frac{q}{2^{256}}
\]
where $q$ is the number of queries and $\sigma$ is the total plaintext length in
128-bit blocks, assuming unique nonces and AES-256 as a PRP.
\end{theorem}

\subsection{Key Management}

The precompile accepts the key as input. Storing keys on-chain defeats the purpose
of encryption (all state is public). The intended usage pattern is:
\begin{enumerate}
  \item Derive the AES key off-chain (from X25519 shared secret, KMS, or MPC).
  \item Pass the key as a transaction parameter (visible to the caller but not stored).
  \item The encrypted output is stored; only parties with the key can decrypt.
\end{enumerate}

For contract-held secrets, the key should be derived via the X25519 precompile
in the same transaction, used for encryption, and never stored.

\subsection{Plaintext Length Limit}

The 65\,536-byte limit on plaintext prevents gas exhaustion attacks. At 6 gas per
byte, the maximum gas for a single encryption is $3\,000 + 6 \times 65\,536 = 396\,216$,
well within a single block's gas limit.

\section{Gas Cost Justification}

AES-NI achieves $\sim$1 cycle per byte on modern CPUs. At 2.45~GHz, encrypting
1~KB takes $\sim$0.4~$\mu$s. The GHASH authentication adds $\sim$0.2 cycles per byte.

\begin{table}[H]
\centering
\begin{tabular}{lrrr}
\toprule
\textbf{Size} & \textbf{Time ($\mu$s)} & \textbf{Gas} & \textbf{Gas/$\mu$s} \\
\midrule
64 B & 0.4 & 3\,384 & --- \\
1 KB & 1.0 & 9\,144 & --- \\
16 KB & 8.2 & 101\,376 & --- \\
64 KB & 32 & 396\,216 & --- \\
\bottomrule
\end{tabular}
\end{table}

The base cost of 3\,000 gas covers key expansion (14 rounds of AES key schedule)
and GCM initialization. The per-byte cost of 6 gas reflects the amortized AES-NI +
GHASH throughput.

\section{Implementation Notes}

The Go implementation wraps \texttt{crypto/aes} and \texttt{crypto/cipher} from
the standard library, which uses AES-NI on amd64 and arm64 (via NEON) when
available. The implementation is constant-time via hardware instructions.

\begin{lstlisting}[language=Go]
func (c *AES256GCM) encrypt(key, nonce, aad, pt []byte) ([]byte, error) {
    block, err := aes.NewCipher(key)
    if err != nil {
        return nil, err
    }
    gcm, err := cipher.NewGCM(block)
    if err != nil {
        return nil, err
    }
    return gcm.Seal(nil, nonce, pt, aad), nil
}
\end{lstlisting}

\section{Conclusion}

The AES-256-GCM precompile brings the standard authenticated encryption primitive
to the EVM at hardware-accelerated speed. The linear gas model accurately reflects
computation cost. Combined with the X25519 or X-Wing KEM precompiles for key
agreement, AES-256-GCM enables practical on-chain encryption workflows.

\begin{thebibliography}{9}
\bibitem{sp800-38d}
  M.~Dworkin, ``Recommendation for Block Cipher Modes of Operation:
  Galois/Counter Mode (GCM) and GMAC,'' NIST SP~800-38D, 2007.

\bibitem{aes-fips}
  NIST, ``Advanced Encryption Standard (AES),'' FIPS~197, 2001 (revised 2023).
\end{thebibliography}

\end{document}
